La arquitectura de la información en la educación superior atraviesa una transición crítica desde silos de datos centralizados hacia ecosistemas de Libros Mayores Distribuidos (DLT). En este paradigma, la integridad de los registros académicos trasciende la mera disponibilidad, situándose en el dominio de la verificabilidad criptográfica de estado. Los sistemas tradicionales de gestión académica, dependientes de bases de datos relacionales y perímetros de seguridad convencionales, presentan vulnerabilidades estructurales ante la manipulación de metadatos y la falta de transparencia en la persistencia de estados \cite{sharples2016blockchain}. El advenimiento de la tecnología blockchain, y específicamente de redes multichain heterogéneas como Polkadot, ofrece una solución basada en el consenso compartido y la finalidad determinista.

Sin embargo, el desarrollo sobre el framework Substrate introduce desafíos de ingeniería de software que no tienen parangón en el desarrollo off-chain. A diferencia de las arquitecturas monolíticas, una parachain de Substrate opera mediante una Función de Transición de Estado (STF) compilada a WebAssembly (WASM), lo que permite actualizaciones forkless pero incrementa exponencialmente la superficie de ataque lógica. En este entorno, la inmutabilidad de los smart contracts y la lógica del runtime actúan como un arma de doble filo: mientras garantizan la resistencia a la censura, convierten cualquier error de desbordamiento aritmético o fallo en el control de acceso en una vulnerabilidad permanente e irreversible \cite{tolmach2021survey}.
El Problema de la Integridad en Sistemas Descentralizados

El problema técnico fundamental reside en la asincronía y el no-determinismo de las interacciones entre cadenas. El uso del formato de mensajes Cross-Consensus Messaging (XCM) para la transferencia de credenciales académicas entre diferentes dominios de consenso introduce riesgos de bloqueo de estado (deadlocks) y fallos en la entrega de mensajes que pueden comprometer la atomicidad de los registros. La literatura actual subraya que el Aseguramiento de la Calidad (QA) en Web3 no puede limitarse a pruebas unitarias aisladas, ya que estas no logran capturar las propiedades emergentes de la interacción entre el runtime, la capa de almacenamiento (Merkle-Patricia Tries) y los mecanismos de incentivos económicos \cite{stefanko2023paraspell}.

Además, la gestión de registros de alta fidelidad, como las actas de notas, requiere una precisión matemática absoluta. Un fallo en la lógica de cálculo de promedios ponderados dentro de una pallet de Rust/ink! podría derivar en una divergencia de estado que, bajo un modelo de Seguridad Compartida, afectaría no solo a la base de datos local, sino a la validez de las pruebas de validez enviadas a la Relay Chain \cite{gao2022blockchain}. Actualmente, existe una brecha metodológica: no hay un marco de trabajo que integre la Verificación Formal (formal methods) y el Mutation Testing específicamente diseñado para los flujos de trabajo de una universidad descentralizada.
Contribución y Objetivos

Para mitigar estos riesgos, esta investigación propone un Framework de QA Multidimensional. El objetivo principal es diseñar un protocolo de validación exhaustivo que garantice la integridad de los datos académicos mediante el uso de métodos formales para la verificación de invariantes lógicas y técnicas de inyección de fallos (mutantes) para evaluar la resiliencia de la suite de pruebas frente a vectores de ataque conocidos. El estudio se desglosa en los siguientes objetivos tácticos:
\begin{itemize}
    \item Modelar formalmente la lógica de transición de estado de una pallet de gestión académica para eliminar ambigüedades en el cálculo de créditos.

    \item Desarrollar una infraestructura de simulación para mensajes XCM que valide la consistencia de los datos en escenarios de red adversos.

    \item Establecer métricas de puntuación de mutación para cuantificar la capacidad de detección de la suite de pruebas frente a errores de lógica de negocio.
\end{itemize}

Este documento se estructura de la siguiente manera: la Sección II presenta una revisión exhaustiva de los trabajos relacionados y el Estado del Arte